%% COMP 9060 Group Project -- Research Proposal (Week 7/8, 15%)
%% Option B: Modern Algorithm Application -- FlashAttention-4
%% proposal_draft.tex -- SKELETON ONLY
%%   Top-level sections = Research Proposal rubric criteria.
%%   Subsections        = required Option B proposal content under each criterion.
%% Format: ACM Conference Proceedings Primary Article Template (acmart, sigconf)
\documentclass[sigconf,nonacm]{acmart}

\settopmatter{printacmref=false}
\renewcommand\footnotetextcopyrightpermission[1]{}
\pagestyle{plain}

%% Italic run-in lead-in with normal paragraph spacing (acmart-safe; avoids the
%% large vertical skip that \paragraph inserts).
\newcommand{\lead}[1]{\textit{#1}\quad}

%% Modest spacing between paragraphs for readability (acmart-safe; does not
%% redefine sectioning commands).
\setlength{\parskip}{0.4\baselineskip}

%% For the tiling / online-softmax schematic in Appendix C.
\usepackage{tikz}
\usetikzlibrary{arrows.meta,positioning,decorations.pathreplacing}

%% For algorithm pseudocode blocks in the appendices.
\usepackage{algorithm}
\usepackage{algpseudocode}

\begin{document}

\title{A Reproducibility Benchmark of the FlashAttention Family Against Classical and Approximate Attention}

\author{Bahaa Al-Juboori}
\affiliation{%
  \institution{British Columbia Institute of Technology}
  \city{Burnaby}
  \state{BC}
  \country{Canada}}

\author{Shaaz E Jahan}
\affiliation{%
  \institution{British Columbia Institute of Technology}
  \city{Burnaby}
  \state{BC}
  \country{Canada}}

\author{Jason Fakidis}
\affiliation{%
  \institution{British Columbia Institute of Technology}
  \city{Burnaby}
  \state{BC}
  \country{Canada}}

\author{Aneet Kaur}
\affiliation{%
  \institution{British Columbia Institute of Technology}
  \city{Burnaby}
  \state{BC}
  \country{Canada}}

\renewcommand{\shortauthors}{Al-Juboori, Jahan, Fakidis, Kaur}

\begin{abstract}
\itshape
Transformer self-attention costs time and memory that grow quadratically in
sequence length, and on modern GPUs this cost is dominated by data movement
rather than arithmetic. The FlashAttention family addresses this with IO-aware
exact attention: it reorganizes the computation to respect the GPU memory
hierarchy, achieving large speedups and linear memory with no approximation. We
propose to benchmark this family, from the original FlashAttention through its
hardware-specialized successors, against traditional exact and approximate
baselines. Using single-GPU micro-benchmarks that sweep sequence length, masking,
and pass type, we measure runtime, throughput, peak memory, maximum context
before out-of-memory, and numerical accuracy. Our goal is to test whether the
claimed advantages hold on accessible hardware, to isolate how much of the
advantage comes specifically from IO-awareness, and to document the practical
limits of adoption, including the alpha-stage status of the newest kernel and the
degree to which production libraries have already absorbed its techniques.
\end{abstract}

\keywords{attention, Transformers, GPU kernels, IO-awareness, FlashAttention}

\maketitle

%% ===== Rubric criterion 1: Problem Definition & Motivation (25%) =====
\section{Problem Definition \& Motivation}
\label{sec:problem}

\subsection{Problem}
Self-attention is the mechanism that lets a Transformer~\cite{vaswani2017attention}
relate every position in a sequence to every other. For a sequence of length $N$,
each position is represented by a query, key, and value vector of dimension $d$,
stacked into matrices $Q,K,V \in \mathbb{R}^{N\times d}$. Attention computes a
score between every query and every key, normalizes these scores with a row-wise
softmax, and uses them to take a weighted average of the values:
$\mathrm{softmax}(QK^\top/\sqrt{d})\,V$. The intermediate score matrix $QK^\top$
therefore has one entry for every pair of positions.

That pairwise structure is the source of the cost. The score matrix is
$N\times N$, so forming it, applying softmax, and multiplying by $V$ take time and
memory that grow quadratically in $N$. On a modern GPU the binding constraint is
not the arithmetic but the \emph{memory traffic}: the processor's compute
throughput far exceeds its memory bandwidth, and the memory hierarchy is steeply
tiered, with a small, fast on-chip SRAM sitting above a large, slow
high-bandwidth memory (HBM). A naive implementation writes the full $N\times N$
matrix to HBM and reads it back, so attention becomes \emph{memory-bound}: its
speed is governed by moving the quadratic intermediate to and from HBM, not by
the floating-point operations.

The consequence is that long context is expensive precisely where it is most
useful. As $N$ grows, the quadratic memory and traffic quickly dominate runtime
and exhaust GPU memory, which has historically capped the sequence lengths
Transformers can handle. The problem we study is whether this bottleneck can be
removed \emph{without} changing the function computed, by reorganizing the
computation to respect the memory hierarchy rather than approximating attention,
and, concretely, whether the FlashAttention family delivers that in practice and
under what conditions.

\subsection{Motivation}
Long context has become central to modern machine learning: large language
models, long-document understanding, and high-resolution vision all depend on
attending over many thousands of positions, and the quadratic cost of attention
is the binding constraint on how far they can scale. Making exact attention
faster and more memory-efficient therefore pays off broadly, improving latency,
throughput, and the maximum context that fits on a given GPU. Crucially, the
IO-aware approach achieves this \emph{without} approximating attention, so it
avoids the accuracy penalty that sub-quadratic methods incur; where it applies,
it is close to a free improvement, which is why it has largely displaced earlier
efficiency techniques in practice.

This makes the FlashAttention family worth studying empirically rather than
taking its advantages for granted. Its performance claims are reported by the
authors on specific, often high-end, hardware, while the production libraries
(cuDNN and PyTorch SDPA) have progressively absorbed the same techniques, so how
much advantage remains in practice is an open question. The newest version,
FlashAttention-4, is both very recent and an alpha-stage release, which raises
real questions of software maturity, numerical precision, and reproducibility. An
independent benchmark on accessible hardware, spanning the lineage from the
original algorithm to its latest variant, is therefore a useful contribution: it
tests whether the claimed benefits hold, isolates where they come from, and
documents the practical cost of adopting them.

%% ===== Rubric criterion 2: Literature Review (35%) =====
\section{Literature Review}
\label{sec:litreview}
We review the lines of research that led to our modern algorithm, following the
intellectual lineage that FlashAttention~\cite{dao2022flashattention} itself
traces. Each paragraph treats one line of work; a theoretical comparison and the
supporting mathematics are given in the supplementary appendices.

\lead{Origins of attention and the quadratic problem.}
Attention originated as a memory-programming idea: Schmidhuber's fast-weight
controllers, in which a slow network learns to generate the weights of a fast
network through outer products~\cite{schmidhuber1992fastweights}, were later
reframed explicitly as ``attention to the recent past''~\cite{ba2016fastweights}.
The modern, content-based form was introduced for neural machine
translation~\cite{bahdanau2015nmt} and then made the center of an architecture by
the Transformer~\cite{vaswani2017attention}. The Transformer's self-attention,
however, has time and memory cost quadratic in sequence length $N$, which is the
bottleneck all later work addresses.

\lead{IO-aware runtime optimization.}
A separate, older line from systems and computer architecture studies how to
optimize reads and writes across a memory hierarchy. The input/output complexity
model~\cite{aggarwal1988io} formalizes counting accesses to slow versus fast
memory, and the Roofline model~\cite{williams2009roofline} frames performance in
terms of arithmetic intensity against memory bandwidth. Applied to Transformers,
this perspective shows that data movement, not arithmetic, is the dominant
cost~\cite{ivanov2021datamovement}. This line supplies the key principle that our
modern algorithm brings to attention.

\lead{Approximate (sub-quadratic) attention.}
The first response to the quadratic cost was to approximate attention. Sparse and
hashing-based methods such as Reformer reduce cost by attending to a
subset~\cite{kitaev2020reformer}, while low-rank and kernel methods achieve
linear complexity: Performer uses random features~\cite{choromanski2021performer},
Linformer projects to a low-rank subspace~\cite{wang2020linformer}, and linear
attention reformulates the computation as a recurrence~\cite{katharopoulos2020linear}.
These methods reduce FLOPs but often fail to deliver wall-clock speedup, because
theoretical efficiency does not automatically translate to hardware
efficiency, a mismatch known as the hardware
lottery~\cite{hooker2020hardwarelottery}. This gap is precisely what motivates an
IO-aware, exact approach.

\lead{Exact, memory-reduced attention.}
A second response keeps attention exact and attacks memory instead. The online
softmax normalizer computes the softmax incrementally over blocks without
materializing a full row~\cite{milakov2018online}, the primitive that makes
tiled exact attention possible. Rabe and Staats then showed self-attention needs
only sub-quadratic \emph{memory} while remaining exact, though its time stays
$O(N^2)$~\cite{rabe2021selfattention}. These exact-but-not-IO-aware methods are
the natural comparison points for our modern algorithm.

\lead{Synthesis: the FlashAttention lineage.}
FlashAttention~\cite{dao2022flashattention} unifies IO-awareness with exact
attention, using tiling and online softmax to keep the $N\times N$ matrix out of
high-bandwidth memory. The result is exact (not an approximation), faster in
wall-clock time, and linear in memory. Its stated limitation, that each kernel
must be hand-written in CUDA, drove the subsequent hardware co-design line:
FlashAttention-2~\cite{dao2023flashattention2},
FlashAttention-3~\cite{shah2024flashattention3}, and our modern algorithm
FlashAttention-4~\cite{zadouri2026flashattention4}, which this project
benchmarks.

%% ===== Option B required content: Modern Algorithm & Baseline Selection =====
\section{Methods Under Study}
\label{sec:methods}
This section identifies the modern algorithm we study and the baselines we
compare it against: Section~\ref{subsec:modern} selects the FlashAttention family
and explains why it is an advancement, and Section~\ref{subsec:baselines} lists
the traditional methods used for comparison.

\subsection{Modern Algorithm Selection}
\label{subsec:modern}
\lead{Choice of modern algorithm.}
We select the FlashAttention family of IO-aware exact-attention algorithms as our
modern algorithm, centering our study on the original FlashAttention
(FA1)~\cite{dao2022flashattention}, introduced in 2022. FA1 established the core
principle the whole family shares: computing exact attention while minimizing
data movement between GPU high-bandwidth memory and on-chip SRAM, using tiling
together with an online-softmax formulation. We treat FA1 as the primary object
of study because it defines this idea, and we additionally examine its modern
variants: FlashAttention-2 (2023)~\cite{dao2023flashattention2},
FlashAttention-3 (2024)~\cite{shah2024flashattention3}, and FlashAttention-4
(2026)~\cite{zadouri2026flashattention4}, which keep FA1's algorithmic core
unchanged while re-engineering the kernel for successive GPU generations.

\lead{Why it is an advancement.}
FA1's central advance is a change of perspective: because attention is
memory-bound rather than compute-bound on GPUs, the cost that matters is data
movement between high-bandwidth memory (HBM) and on-chip SRAM, not arithmetic.
By tiling the computation and using an online-softmax formulation, FA1 avoids
ever writing the $N\times N$ score matrix to HBM, which reduces memory from
quadratic to linear in sequence length and cuts HBM traffic
accordingly~\cite{dao2022flashattention}. Crucially this is achieved with
\emph{no approximation}: FA1 computes exactly the same function as standard
attention, so it delivers wall-clock speedups (up to $7.6\times$ on GPT-2 in the
original paper) without the accuracy trade-off that the approximate methods in
Section~\ref{sec:litreview} incur, and it enables context lengths that were
previously infeasible. Each later version keeps this exact IO-aware core and
targets the next hardware bottleneck:
FlashAttention-2~\cite{dao2023flashattention2} improves work partitioning across
thread blocks and warps and reduces non-matmul operations, roughly doubling
throughput on A100; FlashAttention-3~\cite{shah2024flashattention3} exploits
Hopper asynchrony (warp specialization and overlapped matmul/softmax) and FP8
low precision; and FlashAttention-4~\cite{zadouri2026flashattention4} co-designs
the algorithm and kernel for Blackwell, overlapping matrix multiplication with
the softmax and memory work that now dominate.

\lead{Application.}
We apply the FlashAttention family to long-sequence Transformer attention,
studied as a controlled single-GPU micro-benchmark rather than full model
training: we isolate one attention layer and sweep sequence length, masking, and
pass type, measuring runtime, throughput, peak memory, and the maximum context
length before out-of-memory. This application lets us
test, on accessible GPU hardware, whether FA1's IO-aware principle and its
hardware-specialized successors deliver their claimed benefits in practice.

\subsection{Baseline Selection}
\label{subsec:baselines}
\lead{Baselines.}
We compare FlashAttention against four traditional baselines that predate it and
span the two established families of attention computation, exact and
approximate, exceeding the three classical comparisons required for the Excellent
band. These are
the methods against which FlashAttention's original contribution is properly
measured; the later FlashAttention versions and the production cuDNN/SDPA kernel
are the same algorithmic family and are therefore used as references in the
lineage study rather than as traditional baselines.

\lead{Naive scaled dot-product attention.}
The standard implementation~\cite{vaswani2017attention} materializes the full
$N\times N$ score matrix in GPU memory, giving $O(N^2)$ memory and memory traffic.
It is the canonical ``before FlashAttention'' reference and the method whose
quadratic cost motivated the entire line of work.

\lead{Rabe-Staats memory-efficient attention.}
This method~\cite{rabe2021selfattention} computes \emph{exact} attention with
sub-quadratic memory via lazy-softmax chunking, but its time stays $O(N^2)$ and
it is not IO-aware. Including it lets us isolate how much of FlashAttention's
advantage comes specifically from IO-aware tiling rather than chunking alone.

\lead{Approximate attention (Performer, Linformer).}
These represent the sub-quadratic family that trades exactness for scalability:
Performer approximates the softmax kernel with random
features~\cite{choromanski2021performer}, and Linformer projects keys and values
to a low-rank subspace~\cite{wang2020linformer}. Both achieve near-linear cost
but introduce an accuracy gap relative to exact attention, the trade-off
FlashAttention avoids, and the basis of our exact-versus-approximate comparison.

\lead{PyTorch SDPA / cuDNN (production reference).}
We additionally report the fused SDPA/cuDNN backend as a strong production
reference. Because modern SDPA dispatches to FlashAttention-class kernels, it is
not a traditional baseline; we include it to show how a well-engineered library
compares on each GPU, and to avoid overstating gains against a naive strawman.

%% ===== Rubric criterion 3: Plan & Methodology (30%) =====
\section{Plan \& Methodology}
\label{sec:plan}
This section states the questions the study will answer and the experiments
designed to answer them: Section~\ref{subsec:rqs} sets out the research
questions, and Section~\ref{subsec:expplan} describes the benchmark design,
swept parameters, metrics, and reproducibility measures.

\subsection{Research Questions}
\label{subsec:rqs}
Our research questions address the advantages claimed, the conditions under which
we expect them, and the potential limitations; the experiments that will test
them are described in the Experimental Plan below.

\lead{Claimed advantages.}
Does FlashAttention's IO-aware exact attention deliver its claimed advantages
over traditional baselines, namely lower runtime, higher throughput, and $O(N)$
rather than $O(N^2)$ memory (RQ1)? We further ask how much of this advantage comes
specifically from IO-aware tiling rather than chunking alone, isolated by
comparison with the exact but non-IO-aware Rabe-Staats
method~\cite{rabe2021selfattention} (RQ2); and whether FlashAttention matches the
efficiency of approximate methods such as
Performer~\cite{choromanski2021performer} and
Linformer~\cite{wang2020linformer} while incurring no approximation error (RQ3).

\lead{Conditions for better performance.}
Under what conditions is the advantage largest? We expect it to grow with
sequence length, where the quadratic cost dominates, and to differ between causal
and non-causal masking, since causal masking lets the kernel skip roughly half of
the score blocks (RQ5). We also ask whether the advantage holds across GPU
generations: does each FlashAttention version reproduce its claimed improvement
over its predecessor, with FlashAttention-2 over
FlashAttention-1~\cite{dao2023flashattention2} on Ampere,
FlashAttention-3~\cite{shah2024flashattention3} on Hopper, and
FlashAttention-4~\cite{zadouri2026flashattention4} on Blackwell (RQ4)?

\lead{Potential limitations.}
What are the practical limits of adopting FlashAttention (RQ6)? These include
hardware dependence (each version targets a specific GPU generation), software
maturity (the FlashAttention-4 kernel is an alpha-stage release), numerical
precision across BF16/FP8/FP4, and reproducibility. We further ask whether the
advantage is durable: because the production cuDNN/SDPA library has progressively
absorbed FlashAttention's techniques, with FlashAttention-4 reporting only a
modest margin over recent cuDNN, how much benefit remains for the research kernel
on current hardware (RQ7)?

\subsection{Experimental Plan}
\label{subsec:expplan}
We evaluate FlashAttention through controlled single-GPU micro-benchmarks of one
attention layer. The plan comprises a primary study against traditional baselines
(RQ1 to RQ3), a per-generation lineage study (RQ4), and shared limitation analyses
(RQ6 and RQ7).

\lead{Methods.}
The methods under study are FlashAttention-1/2,
FlashAttention-3 (Hopper build), and FlashAttention-4. Traditional baselines are
naive attention~\cite{vaswani2017attention},
Rabe-Staats memory-efficient attention~\cite{rabe2021selfattention}, and the
approximate methods Performer~\cite{choromanski2021performer} and
Linformer~\cite{wang2020linformer}; PyTorch SDPA/cuDNN serves as a production
reference.

\lead{Swept parameters.}
We sweep three parameters directly: sequence length $N \in \{512, 1\mathrm{K},
2\mathrm{K}, 4\mathrm{K}, 8\mathrm{K}, 16\mathrm{K}, 32\mathrm{K}, 64\mathrm{K}\}$
(extended until out-of-memory); masking (causal vs.\ non-causal), since causal
masking skips roughly half the score blocks; and pass type (forward vs.\
forward+backward), since the backward pass recomputes attention tiles. Head
dimension is fixed at $128$ (with $64$ reported as a secondary point) and batch
size is held constant, so that $N$ is isolated as the primary axis.

\lead{Platform axes.}
Two further axes are coupled to hardware rather than swept freely. \emph{GPU
generation}: FA1/FA2 and the traditional baselines run on Ampere (A100), FA3 on
Hopper (H100), and FA4 on Blackwell (RTX PRO 6000 Blackwell). \emph{Precision}:
BF16 is used on all platforms for comparability, with FP8 added on Hopper (FA3)
and FP4/MXFP8 on Blackwell (FA4), since these low-precision modes exist only on
the newer hardware. We therefore report BF16 as the common axis and treat FP8/FP4
as platform-specific extensions. These are single-GPU spot-rented instances; the
full sweep requires only a few GPU-hours per platform.

To keep comparisons fair, each comparison is made on a single machine. The
primary study (RQ1 to RQ3), which compares FlashAttention against the traditional
baselines, runs entirely on one Ampere (A100) machine, with every method in the
same language and under the same BF16 precision; we report that machine's
hardware and software. The lineage study (RQ4) necessarily spans machines because
FA3 and FA4 require Hopper and Blackwell respectively and cannot run elsewhere; we
therefore frame it as a set of per-generation comparisons (each FlashAttention
version against its predecessor on the hardware that version targets) rather than
a single cross-machine ranking, and we report the hardware and software of every
machine used.

\lead{Metrics and protocol.}
For each configuration we record forward and forward+backward latency, throughput
(tokens/s and TFLOP/s), peak GPU memory, maximum $N$ before out-of-memory,
and, for approximate methods and low-precision kernels, numerical error
relative to exact attention. All compared methods are implemented in the same
language (Python/PyTorch) and run on the same machine per platform, whose
hardware and software we report. Each configuration is warmed up and repeated at
least ten times per input; we report the \emph{median} running time and its
spread (minimum--maximum, or interquartile range), and we plot running time
against sequence length for each method. Correctness is verified by asserting
that each exact method agrees with naive attention within floating-point
tolerance.

\lead{Reproducibility.}
We pin and report GPU model, driver, CUDA, PyTorch, and kernel versions, and
release all benchmark scripts. The main practical risks, the alpha-stage
FlashAttention-4 kernel and precision-dependent numerical behavior, are
documented as part of the limitations analysis (RQ6 and RQ7).

%% ===== Rubric criterion 4: Team Organization (10%) =====
\section{Team Organization}
\label{sec:team}

\subsection{Role Assignments}
Table~\ref{tab:roles} assigns each member a primary area of responsibility that
maps onto the structure of the study. All members contribute to the final paper
and presentation.

\begin{table}[h]
  \centering
  \caption{Primary role assignments.}
  \label{tab:roles}
  \small
  \begin{tabular}{@{}lp{0.62\linewidth}@{}}
    \toprule
    \textbf{Member} & \textbf{Primary responsibility} \\
    \midrule
    Shaaz E Jahan & Mathematical and algorithmic preliminaries for all methods: the
      naive implementation, each FlashAttention variant, and the baselines. \\
    Jason Fakidis & The primary study: experiments comparing FlashAttention-1
      against the traditional baselines. \\
    Bahaa Al-Juboori & The lineage study: FA1/2/3/4 comparisons across GPU
      generations (Ampere, Hopper, Blackwell). \\
    Aneet Kaur & Validation and discussion: verifying each experiment meets the
      outcomes set by the research questions, interpreting the results, and
      preparing the discussion. \\
    \bottomrule
  \end{tabular}
\end{table}

\subsection{Timeline}
Table~\ref{tab:timeline} gives the week-by-week plan for Weeks~8 to~14. Work is
sequenced so that the preliminaries precede implementation, the two experimental
studies run in parallel once the environment is ready, and validation follows the
collection of results.

\begin{table}[h]
  \centering
  \caption{Project timeline (Weeks 8 to 14).}
  \label{tab:timeline}
  \small
  \begin{tabular}{@{}lp{0.66\linewidth}@{}}
    \toprule
    \textbf{Week} & \textbf{Activity} \\
    \midrule
    8 & Finalize scope, methods, and research questions; stand up the shared
      repository; secure GPU access. Shaaz begins the mathematical and
      algorithmic preliminaries. \\
    9 & Preliminaries completed; implement and verify the baselines and the
      benchmark harness (timing, memory, out-of-memory capture). \\
    10 & Pilot runs: confirm FA1/FA2 and the baselines run correctly on Ampere,
      and confirm FA3 (Hopper) and FA4 (Blackwell) compatibility. \\
    11 & Jason runs the primary study (FA1 vs.\ baselines); Bahaa begins the
      lineage study (FA1/2/3/4 across Ampere, Hopper, Blackwell). \\
    12 & Complete both studies; collect at least five repeated measurements per
      input and record the median and spread. \\
    13 & Aneet validates results against the research-question outcomes;
      interpret findings and draft the discussion; produce figures and tables.
      Prepare and rehearse the presentation (delivered this week); submit slides
      as a PDF. \\
    14 & Finalize the paper and repository documentation; submit the contribution
      statement and peer evaluation. \\
    \bottomrule
  \end{tabular}
\end{table}

\bibliographystyle{ACM-Reference-Format}
\bibliography{references}

\clearpage
\appendix

\section{Requirements Coverage and Rubric Mapping}
\label{app:coverage}
This appendix lists every deliverable requirement and rubric item for the
research proposal and maps each to where this document satisfies it. Table~
\ref{tab:content} covers the Option~B required-content items; Table~
\ref{tab:rubric} covers the grading rubric; the remaining items cover format and
submission requirements.

\begin{table}[h]
  \centering
  \caption{Option B required-content items and where each is addressed.}
  \label{tab:content}
  \small
  \begin{tabular}{@{}p{0.46\linewidth}p{0.44\linewidth}@{}}
    \toprule
    \textbf{Required content} & \textbf{Where addressed} \\
    \midrule
    Problem and motivation &
      Sec.~\ref{sec:problem} (Problem and Motivation) \\
    Modern algorithm: which, cited, with year &
      Sec.~\ref{subsec:modern}, ``Choice of modern algorithm'' (FlashAttention,
      2022, with the 2023/2024/2026 variants) \\
    Modern algorithm: why an advancement &
      Sec.~\ref{subsec:modern}, ``Why it is an advancement'' \\
    Modern algorithm: problem/application &
      Sec.~\ref{subsec:modern}, ``Application'' \\
    Baseline selection (2--4 traditional) &
      Sec.~\ref{subsec:baselines} (naive, Rabe-Staats, Performer, Linformer;
      SDPA/cuDNN as reference) \\
    Research questions: claimed advantages &
      Sec.~\ref{subsec:rqs}, ``Claimed advantages'' (RQ1--RQ3) \\
    Research questions: conditions &
      Sec.~\ref{subsec:rqs}, ``Conditions'' (RQ4--RQ5) \\
    Research questions: limitations &
      Sec.~\ref{subsec:rqs}, ``Potential limitations'' (RQ6--RQ7) \\
    Research questions: experiments &
      Sec.~\ref{subsec:expplan} (Experimental Plan) \\
    Team plan: roles and timeline &
      Sec.~\ref{sec:team} (Role Assignments and Timeline) \\
    \bottomrule
  \end{tabular}
\end{table}

\begin{table}[h]
  \centering
  \caption{Research Proposal rubric (15\% of project grade) and how this
  document addresses each criterion.}
  \label{tab:rubric}
  \small
  \begin{tabular}{@{}p{0.22\linewidth}cp{0.50\linewidth}@{}}
    \toprule
    \textbf{Criterion} & \textbf{Wt.} & \textbf{How addressed} \\
    \midrule
    Problem Definition \& Motivation & 25\% &
      Sec.~\ref{sec:problem}: first-principles problem statement and two-paragraph
      motivation \\
    Literature Review & 35\% &
      Sec.~\ref{sec:litreview}: five categories tracing the lineage, with 10+
      cited sources and a theory table (in the supplementary appendices) \\
    Plan \& Methodology & 30\% &
      Sec.~\ref{sec:plan}: research questions (RQ1--RQ7) and a detailed
      experimental plan with swept parameters, metrics, and reproducibility \\
    Team Organization & 10\% &
      Sec.~\ref{sec:team}: role assignments and a week-by-week timeline \\
    \bottomrule
  \end{tabular}
\end{table}

\noindent\textbf{Format and submission requirements.}
\begin{itemize}
  \item \emph{ACM template}: the document uses the ACM Primary Article
    (\texttt{acmart}, \texttt{sigconf}) template.
  \item \emph{Length (1--2 pages)}: the main body targets the required length;
    appendices and references are supplementary.
  \item \emph{Citations}: all sources are cited in ACM reference format via
    BibTeX.
\end{itemize}

\end{document}
